% GENERATED FILE. Edit canonical lessons, then run npm run generate:books.
% canonical-source-hash: fnv1a64:9730e9978c149632
% canonical-lessons: KA-C09-kshamisi, KA-S118-letter-dda, KA-S140-digit-one

\chapter{Sorry}
\label{ch:ka-sorry}
\hlchaptermodality{9}

\begin{chapteropening}
\textbf{By the end of this chapter:} \emph{I can apologise in Kannada and ask someone to forgive me.}

Say \kn{ಕ್ಷಮಿಸಿ}, trace it back through kṣamisu to Sanskrit kṣamā, and contrast it with Tamil's unrelated maṉṉikkavum.
\end{chapteropening}

\section[\texorpdfstring{\kn{ಕ್ಷಮಿಸಿ} (kṣamisi)}{ಕ್ಷಮಿಸಿ (kṣamisi)}]{\kn{ಕ್ಷಮಿಸಿ} (kṣamisi) — forgive me / sorry}

\label{lesson:KA-C09-kshamisi}



\begin{quote}
You already know \textbf{\kn{ದಯವಿಟ್ಟು}} (\emph{dayaviṭṭu}), \textquotedblleft{}please.\textquotedblright{} Kannada's \textquotedblleft{}sorry\textquotedblright{} reuses a trick you'll see everywhere in this language: taking a Sanskrit noun and turning it into a working Kannada verb.
\end{quote}

\begin{cousinweb}
\begin{itemize}
\raggedright
  \item \textbf{\kn{ಕ್ಷಮಾ}} (\emph{kṣamā}) = \textquotedblleft{}\textbf{patience, forbearance, forgiveness}\textquotedblright{} — a Sanskrit noun, unchanged.
  \item \textbf{‑\kn{ಇಸು}} (\emph{‑isu}) = a \textbf{verb-making suffix}: bolt it onto a Sanskrit noun and you get a brand-new Kannada verb. \emph{kṣamā} + \emph{‑isu} $\to$ \textbf{\kn{ಕ್ಷಮಿಸು}} (\emph{kṣamisu}), \textquotedblleft{}\textbf{to forgive}.\textquotedblright{} (The same move builds \emph{prārambhisu}, \textquotedblleft{}to begin,\textquotedblright{} from Sanskrit \emph{prārambha}.)
  \item \textbf{‑\kn{ಿ}} (\emph{‑i}) = the plain \textbf{imperative} ending — the very one you met on \emph{kuḷitukoḷḷi} (\textquotedblleft{}please sit\textquotedblright{}) in the please lesson.
\end{itemize}

So \textbf{\kn{ಕ್ಷಮಿಸಿ}} is simply \textquotedblleft{}\textbf{forgive!}\textquotedblright{} — \emph{kṣamisu} wearing its imperative ending.
\end{cousinweb}

\begin{cousinweb}
Three of the four Dravidian cousins build \textquotedblleft{}sorry\textquotedblright{} from the very same Sanskrit noun, each with its own verb-making suffix and its own imperative shape:

\begin{itemize}
\raggedright
  \item Kannada \textbf{\kn{ಕ್ಷಮಿಸಿ}} \emph{kṣami\textbf{si}} — Sanskrit \emph{kṣamā} + \textbf{‑isu} (verb) + \textbf{‑i} (imperative)
  \item Telugu \textbf{\te{క్షమించండి}} \emph{kṣamin̄\textbf{caṇḍi}} — Sanskrit \emph{kṣamā} + \textbf{‑in̄cu} (verb) + \textbf{‑aṇḍi} (imperative)
  \item Malayalam \textbf{\ml{ക്ഷമിക്കണം}} \emph{kṣami\textbf{kkaṇaṁ}} — Sanskrit \emph{kṣama} + \textbf{‑ikkuka} (verb) + \textbf{‑aṇaṁ} (necessitative)
  \item Tamil \textbf{\ta{மன்னிக்கவும்}} \emph{maṉṉi\textbf{kkavum}} — its \textbf{own} native root \emph{maṉṉi}, not Sanskrit at all
\end{itemize}
\end{cousinweb}

\begin{culture}
\textbf{\kn{ಕ್ಷಮಿಸಿ}} is a genuine, everyday way to say \textquotedblleft{}sorry\textquotedblright{} — not overly formal — but in casual conversation you'll also hear the softer \textbf{\kn{ಕ್ಷಮಿಸಿ} \kn{ಬಿಡಿ}} (\emph{kṣamisi biḍi}, roughly \textquotedblleft{}just forgive {[}it{]}\textquotedblright{}) or simply the borrowed English \textquotedblleft{}\textbf{sorry},\textquotedblright{} which is extremely common in spoken Kannada.
\end{culture}

\subsection*{Your turn}
\begin{itemize}
\raggedright
  \item \emph{Say it:} \textquotedblleft{}kṣamā\textquotedblright{} — the Sanskrit noun, patience/forgiveness
  \item \emph{Say it:} the verb — \textquotedblleft{}kṣamisu,\textquotedblright{} to forgive — then \textquotedblleft{}kṣamisi,\textquotedblright{} forgive!
  \item \emph{Say it:} contrast — Tamil's maṉṉikkavum uses its own root instead
  \item \emph{Recall:} say \emph{aidu}, then say \emph{āru}
\end{itemize}

\begin{tcolorbox}[breakable,colback=teal!4,colframe=teal!35!black,title={Before you move on}]
What is \textbf{\kn{ಕ್ಷಮಿಸಿ}} built from? (\textbf{Sanskrit kṣamā} \textquotedblleft{}forgiveness\textquotedblright{} + the verb-making suffix \textbf{‑isu} + the imperative \textbf{‑i}.) Which cousin languages do the same trick, and which one doesn't? (\textbf{Telugu and Malayalam} also verb-ify \emph{kṣamā}; \textbf{Tamil} uses its own native root \emph{maṉṉi} instead.) What do people often say casually instead? (\textbf{Kṣamisi biḍi}, or simply borrowed English \textquotedblleft{}\textbf{sorry}.\textquotedblright{})
\end{tcolorbox}
\section[\texorpdfstring{\kn{ಡ} (ḍa)}{ಡ (ḍa)}]{\kn{ಡ} — one character, met inside words you already say}

\label{lesson:KA-S118-letter-dda}



\begin{quote}
Before the new one: \kn{ಯ} — what does it do?

One character this time. Just one — and you have been saying it for pages without knowing which mark on the page it was.
\end{quote}

\subsection*{Script you'll notice: \kn{ಡ}}
\textbf{\kn{ಡ}} — \emph{ḍa}.

It is a \textbf{consonant}, and in this script a consonant is never bare: it comes with an \emph{a} already in it. So it is not \emph{ḍ}, it is \textbf{ḍa}.

You already say these, and every one of them has \kn{ಡ} somewhere inside it:

\begin{itemize}
\raggedright
  \item \textbf{\kn{ಮಾತನಾಡು}} \emph{mātanāḍu} — to speak
  \item \textbf{\kn{ಕೆಲಸ} \kn{ಮಾಡು}} \emph{kelasa māḍu} — to work (lit. \textquotedblleft{}work-do\textquotedblright{})
  \item \textbf{\kn{ಎರಡು}} \emph{eraḍu} — two
\end{itemize}

\subsection*{Writing: \kn{ಡ}}
\hlblockfigure{\detokenize{figures/KA-S118-letter-dda-filmstrip.pdf}}{How \kn{ಡ} is written, stroke by stroke}

\begin{itemize}
\raggedright
  \item \textbf{1.} start at the upper left and go down the left side into the left lobe, rising into the middle point
  \item \textbf{2.} without lifting, drop from the point around the right lobe and climb the right side
  \item \textbf{3.} without lifting, curl left into the small inner loop and go round it counterclockwise
  \item \textbf{4.} without lifting, rise out of the loop and close the bowl leftward along the top
  \item \textbf{5.} lift, then draw the top bar from left to right
  \item \textbf{6.} without lifting, curl up into the hook
\end{itemize}

\textbf{Pen lifts: 1.} The whole body, with the point in its base and the small loop inside its right lobe, is one run that closes along the top; the curled bar comes after the only lift. It is \kn{ದ} with that small loop tucked inside.

\begin{quote}
This is one attested teaching order and not a national standard — Kannada handwriting is taught with school-to-school variation. Source: Gopala Krishna A, 'Kannada-alphabet-da.gif', 49 frames, Wikimedia Commons, 25 May 2016, CC BY-SA 4.0.
\end{quote}

\begin{quote}
The animation is filed as 'da': that series keeps 'da' for this letter, \kn{ಡ}, and spells dental \kn{ದ} with an \emph{h}, as 'dha'.
\end{quote}

\subsection*{Your turn}
\begin{itemize}
\raggedright
  \item \emph{Look:} at these words, and find \kn{ಡ} in the ones that have it
\end{itemize}

\begin{quote}
\kn{ಮಾತನಾಡು}  ·  \kn{ನಮಸ್ಕಾರ}
\end{quote}

\begin{itemize}
\raggedright
  \item \emph{Trace it:} \kn{ಡ} three times, saying \emph{ḍa} as you finish each one
  \item \emph{Look:} back at any page of this chapter and find \kn{ಡ} once more
\end{itemize}

\begin{tcolorbox}[breakable,colback=teal!4,colframe=teal!35!black,title={Before you move on}]
Which character is this — \kn{ಡ}? What sound does it carry? (\emph{\textbf{ḍa}}.) Name one word you already say that contains it.
\end{tcolorbox}
\section[\texorpdfstring{\kn{೧} (1)}{೧ (1)}]{\kn{೧} — the digit for ondu}

\label{lesson:KA-S140-digit-one}



\begin{quote}
Before the new shape: \textbf{\kn{ಡ}} — what sound does that letter carry?

You have been counting since you learned to count. Now the written form arrives, one shape at a time.
\end{quote}

\subsection*{Script you'll notice: \kn{೧}}
\textbf{\kn{೧}} — \textbf{1}, the digit for \textbf{\kn{ಒಂದು}} (\emph{ondu}), \textquotedblleft{}one\textquotedblright{}.

Kannada writes its numbers with its own ten digits. You will see them on bus boards, house plates and price tags all over Karnataka, sitting beside the Western digits rather than replacing them. Here is the first.

\subsection*{Writing: \kn{೧} — copy what you see}
Put your pen on \kn{೧} and follow its line. Copy the shape in front of you — slowly, and larger than it is printed.

\begin{quote}
This book does not yet tell you \textbf{where to start the character or which way to travel}. That is taught with real variation from school to school, and it is not written down here until it can be written down with a source. Copying what you can see needs no such source.
\end{quote}

\subsection*{Your turn}
\begin{itemize}
\raggedright
  \item \emph{Look:} at the digits you have met, and name each one aloud
\end{itemize}

\begin{quote}
\kn{೧}
\end{quote}

\begin{itemize}
\raggedright
  \item \emph{Say it:} \textquotedblleft{}ondu\textquotedblright{} as you point at \kn{೧}
\end{itemize}

\begin{tcolorbox}[breakable,colback=teal!4,colframe=teal!35!black,title={Before you move on}]
Which number is \textbf{\kn{೧}}? (\textbf{1} — \emph{ondu}, \kn{ಒಂದು}.) Name every digit you have met so far, in order.
\end{tcolorbox}
